\subsection{线性自同态族的情形}

现在假定 S 具有一个向量空间结构，映射 \(r : S \to \text{End}(V)\) 是线性的，且 V 与 S 都是有限维的。

命题 5. 假定条件 (AC) 满足，并设 \(\lambda : S \to k\) 使 \(V^{\lambda}(S) \neq 0\) 。若 \(k\) 的特征为 0，则映射 \(\lambda\) 是线性的。若 \(k\) 的特征为 \(p \neq 0\) ，则存在 \(q\) （\(p\) 的一个幂）整除 \(\dim V^{\lambda}(S)\) ，以及一个齐次多项式函数 \(P : S \to k\) ，次数为 \(q\) ，使得 \(\lambda(s)^{q} = P(s)\) 对所有 \(s \in S\) 成立。

由于 \(V^{\lambda}(S)\) 在 \(r(S)\) 下稳定（第 1 节的引理 1 与公式 (1)），可假定 \(V = V^{\lambda}(S)\) 。设 \(n = \dim V\) 。于是，对 \(s \in S\) ，

\[
\det (X - r (s)) = (X - \lambda (s)) ^ {n}.
\]

另一方面，行列式的展开表明

\[
\det (\mathrm{X} - r (s)) = \mathrm{X} ^ {n} + a _ {1} (s) \mathrm{X} ^ {n - 1} + \dots + a _ {i} (s) \mathrm{X} ^ {n - i} + \dots
\]

其中 \(a_{i}: S \to k\) 是 i 次齐次多项式函数。写 n = qm，其中 q 是 k 的特征指数的一个幂且 \((q, m) = 1\) 。于是 \((\mathrm{X} - \lambda(s))^{n} = (\mathrm{X}^{q} - \lambda(s)^{q})^{m}\) ；从而 \(-m\lambda(s)^{q} = a_{q}(s)\) ，由此即得结果。

命题 6. 假定 \(k\) 是无限的且条件 (AC) 满足。设 \(k'\) 是 \(k\) 的一个扩张。令 \(V' = V \otimes_k k'\) ，\(S' = S \otimes_k k'\) 。设 \(r': S' \to \text{End}(V')\) 是由 \(r\) 经纯量扩张得到的映射。那么

\[
\mathrm{V} ^ {0} (\mathrm{S}) \otimes_ {k} k ^ {\prime} = \mathrm{V} ^ {\prime 0} (\mathrm{S}) = \mathrm{V} ^ {\prime 0} (\mathrm{S} ^ {\prime}).
\]

第一个等式由命题 1 得出。为证第二个等式，可假定 \(\mathrm{V} = \mathrm{V}^0 (\mathrm{S})\) ，于是 \(\mathrm{V}' = \mathrm{V}'^{0}(\mathrm{S})\) 。设 \((s_1,\dots ,s_m)\) 是 \(\mathrm{S}\) 的一个基，\((e_1,\dots ,e_n)\) 是 \(\mathrm{V}\) 的一个基。存在多项式 \(\mathrm{P}_{ij}(\mathrm{X}_1,\dots ,\mathrm{X}_m)\) 使得

\[
r ^ {\prime} (a _ {1} s _ {1} + \dots + a _ {m} s _ {m}) ^ {n} e _ {j} = \sum_ {i = 1} ^ {n} \mathrm{P} _ {i j} (a _ {1}, \ldots , a _ {m}) e _ {i}
\]

其中 \(1 \leq j \leq n\) 且 \(a_{1}, \ldots, a_{m} \in k'\) 。由假设，\(r'(s)^{n} = 0\) 对所有 \(s \in S\) 成立，换言之，\(\mathrm{P}_{ij}(a_{1}, \ldots, a_{m}) = 0\) 对 \(1 \leq i, j \leq n\) 与 \(a_{1}, \ldots, a_{m} \in k\) 成立。由于 k 是无限的，\(P_{ij} = 0\) 。因此，\(r'(S')\) 的每个元素都是幂零的，且 \(V' = V'^{0}(S')\) 。

命题 7. 假定 \(k\) 是无限的且条件 (AC) 满足。设 \(\tilde{\mathrm{S}}\) 是 \(s \in \mathrm{S}\) 中使 \(\mathrm{V}^0(s) = \mathrm{V}^0(\mathrm{S})\) 的元素所成的集合。若 \(s \in \mathrm{S}\) ，用 \(\mathrm{P}(s)\) 表示 \(\mathrm{V} / \mathrm{V}^0(\mathrm{S})\) 上由 \(r(s)\) 定义的自同态的行列式（第 1 节，定理 1 的推论 2 (i)）。

\begin{enumerate}
\def\labelenumi{(\roman{enumi})}
\item
  函数 \(s \mapsto \mathrm{P}(s)\) 是 \(S\) 上的多项式函数。我们有 \(\tilde{\mathrm{S}} = \{s \in S | \mathrm{P}(s) \neq 0\}\) ；这是 \(S\) 在 Zariski 拓扑中的一个开子集（附录 I）。
\item
  \(\tilde{\mathrm{S}}\) 非空，且 \(\mathrm{V}^{+}(s) = \mathrm{V}^{+}(\mathrm{S})\) 对所有 \(s\in \tilde{\mathrm{S}}\) 成立。
\end{enumerate}

\(s \mapsto \mathrm{P}(s)\) 是多项式函数这一事实由 r 的线性性得出。若 \(s \in \mathrm{S}\) ，则 \(\mathrm{V}^{0}(s) \supset \mathrm{V}^{0}(\mathrm{S})\) ，且等号成立当且仅当 \(r(s)\) 定义 \(\mathrm{V}/\mathrm{V}^{0}(\mathrm{S})\) 的一个自同构，由此得 (i)。

现在设 \(k'\) 是 k 的一个代数闭包，并如命题 6 那样引入 \(V', S', r'\) 。我们指出，\(S'\) 满足条件 (AC)，因为多项式恒等式 \((\operatorname{ad} r(s_1))^{2\dim V-1}r(s_2) = 0\) 对 \(s_1, s_2 \in S\) 成立，且可以延拓（第 1 节，注）。应用定理 1，我们得到一个分解

\[
\mathrm{V} ^ {\prime} = \mathrm{V} ^ {0} (\mathrm{S} ^ {\prime}) \oplus \sum_ {i = 1} ^ {m} \mathrm{V} ^ {\lambda_ {i}} (\mathrm{S} ^ {\prime})
\]

其中 \(\lambda_{i} \neq 0\) 对 \(1 \leq i \leq m\) 成立。对 \(1 \leq i \leq m\) ，存在一个多项式函数 \(\mathrm{P}_{i}\) ，在 \(\mathrm{S}'\) 上非零，以及一个整数 \(q_{i}\) ，使得 \(\lambda_{i}^{q_{i}} = \mathrm{P}_{i}\) （命题 5）。由于 \(k\) 是无限的，存在 \(s \in \mathrm{S}\) 使 \((\mathrm{P}_{1} \ldots \mathrm{P}_{m})(s) \neq 0\) ，参看~代数，第四章，§2，第 3 节，命题 9 的推论 2。于是 \(\lambda_{i}(s) \neq 0\) 对所有 \(i\) 成立，故 \(\mathrm{V}^{\prime 0}(\mathrm{S}') = \mathrm{V}^{\prime 0}(s)\) ，从而 \(\mathrm{V}^{0}(\mathrm{S}) = \mathrm{V}^{0}(s)\) （命题 6），这表明 \(\tilde{\mathrm{S}} \neq \emptyset\) 。若 \(s \in \tilde{\mathrm{S}}\) ，则 \(\mathrm{V}^{+}(\mathrm{S})\) 在 \(r(s)\) 下稳定且是 \(\mathrm{V}^{0}(s)\) 在 V 中的余子空间这一事实蕴含 \(\mathrm{V}^{+}(\mathrm{S}) = \mathrm{V}^{+}(s)\) （定理 1 的推论 2）。

